%------------------------------------------------------------------------------%
\begin{question}
  Encontre uma equação para a tangente à elipse
  \[
    \frac{x^2}{4} + y^2 = 2
  \]
  no ponto \((-2, 1)\)
\end{question}

%------------------------------------------------------------------------------%
\begin{resolution}{Solução}
  Encontre uma equação para a tangente à elipse
  \[
    \frac{x^2}{4} + y^2 = 2
  \]
  no ponto \((-2, 1)\)

  \addgraphic{9}{6.5,2.5}{tangente_elipse}
\end{resolution}

%------------------------------------------------------------------------------%
\begin{resolution}{Solução}
  Precisamos notar que a elipse é uma curva de nível da função
  \[
    f(x, y) = \frac{x^2}{4} + y^2
  \]

  \pause
  Calculando as derivadas parciais de $f$ temos
  \[
    \pause
    \D{f}{x}
    \pause
    = \D{}{x}\left(\frac{x^2}{4} + y^2\right)
    \pause
    = \frac{x}{2}
  \]
  \[
    \pause
    \D{f}{y}
    \pause
    = \D{}{y}\left(\frac{x^2}{4} + y^2\right)
    \pause
    = 2y
  \]
\end{resolution}

%------------------------------------------------------------------------------%
\begin{resolution}{Solução}
  Então o gradiente de $f$, no ponto \((-2, 1)\), é
  \pause
  \[
    \nabla f(-2, 1)
    \pause
    \;=\; \Vetor{\frac{x}{2}}{2y}\evalat{(-2, 1)}{}
    \pause
    = \vetor{-1}{2}
  \]
\end{resolution}

%------------------------------------------------------------------------------%
\begin{resolution}{Solução}
  \onslide<+->{A reta tangente é ortogonal ao vetor gradiente}
  \begin{columns}
  \column{0.5\linewidth}
    \begin{align*}
      \onslide<+->{\nabla f(x_0, y_0) \cdot \vetor{x-x_0}{y-y_0} & = 0 \\}
      \onslide<+->{\nabla f(-2, 1)    \cdot \vetor{x-(-2)}{y-1}  & = 0 \\}
      \onslide<+->{\vetor{-1}{2} \cdot \vetor{x+2}{y-1} & = 0}
    \end{align*}
  \column{0.5\linewidth}
  \begin{align*}
    \onslide<+->{-(x+2) + 2(y-1) & = 0 \\}
    \onslide<+->{-x +2y -4 & = 0 \\}
    \onslide<+->{-x +2y    & = 4 \\}
    \onslide<+->{ x        & = 2y-4}
  \end{align*}
  \end{columns}
\end{resolution}

%------------------------------------------------------------------------------%
